\subsection{Coefficient localization, ordinary tensor products, and integral periods}
\label{arch136:localized-tensor}

\paragraph{Scope.}
These calculations test the common-target route to admissible generalized
de Rham theories in cote~136, physical pp.~117,119,121, archive footers
116,118,120, supplied scans
\texttt{ARCHIVE\_136\_scan117.tex},
\texttt{ARCHIVE\_136\_scan119.tex},
\texttt{ARCHIVE\_136\_scan121.tex}.
They use the inherited category \(\mathcal A_{\mathbf Z}\), the standard
model AQ-136-DP-MODEL, and the incoming replacements
AQ-136-FIRST-SKELETON. They do not decide integral connectedness.
No input A, M, or outside preprint is used.

Throughout put \(S=\Gamma_{\mathbf Z}(T)\), \(t=T^{[1]}\), and
distinguish localization at \(t\) from rationalization:
\[
 S_{\mathbf Q}=\mathbf Q[t],
 \qquad S[t^{-1}]=\mathbf Q[t,t^{-1}]=:\mathcal R.
 \tag{136.LT1}
\]
Indeed \(t^n=n!T^{[n]}\); after inverting \(t\), the element
\(T^{[n]}t^{-n}\) is an inverse to \(n!\). Conversely
\(T^{[n]}=t^n/n!\) in the displayed Laurent polynomial ring.
Weights in a localized module are now allowed to be negative.

\paragraph{ARCH-136-LOCALIZED-TENSOR.}
For any \(C,D\in\mathcal A_{\mathbf Z}\), the ordinary tensor
\(C\otimes_S D\), after localization at \(t\), has coefficient
cohomology \(\mathcal R\) in form degree zero and vanishing positive
vertical cohomology, in every simplex. Every fixed form-degree
horizontal row of the localized tensor is acyclic, including
horizontal degree zero. No flatness hypothesis on the original
integral modules is required.

\emph{Proof.}
Localization is exact and commutes with the normalized chain complex
of a simplicial module. Its component of weight \(w\) is the filtered
colimit of the components of weights \(w+k\), with transitions
multiplication by \(t\). The uniform vertical tail condition for
\(C\) therefore gives
\[
 H^q(C_n[t^{-1}])=
 \begin{cases}
 \mathcal R,&q=0,\\
 0,&q>0.
 \end{cases}
 \tag{136.LT2}
\]
The coefficient map is the indicated isomorphism. Similarly the
uniform horizontal tail condition gives an acyclic localized row
for every form degree. The same holds for \(D\).

Every weight-graded \(\mathcal R\)-module \(M\) is canonically
\(\mathcal R\otimes_{\mathbf Q}M^0\), by multiplication:
the unit \(t\) identifies all weight components with the component
of weight zero. Thus all the relevant tensor calculations reduce
to tensor products of complexes of \(\mathbf Q\)-vector spaces.
Vertically, split a complex with the cohomology (136.LT2) into
its coefficient copy and contractible summands. Its tensor with
the other complex has just the coefficient cohomology.
Horizontally, an acyclic normalized complex of vector spaces
is contractible. The Dold--Kan correspondence turns such a
contraction into a simplicial linear homotopy; tensoring a
simplicial homotopy with another simplicial module, using its
degeneracies in the other factor, still gives a homotopy.
Thus each diagonal tensor row is acyclic. The form-degree-\(q\)
row is the finite sum of these tensors for \(i+j=q\).
Finally localization commutes with tensor products, so these
are precisely the rows of
\((C\otimes_SD)[t^{-1}]\).
\hfill\(\square\)

The conclusion says that the corresponding defects before
localization are individually killed by some power of \(t\).
It gives no uniform exponent and no replacement of \(t\)-power
torsion by integer torsion.

\paragraph{ARCH-136-FLAT-RATIONAL-TENSOR.}
A sufficient hypothesis for the ordinary rational tensor itself
to be admissible is that every module
\((C_n^{*,q})_{\mathbf Q}\) and
\((D_n^{*,q})_{\mathbf Q}\) be flat over \(\mathbf Q[t]\).
For this statement assume \(C_0=D_0=S\), as supplied by
FIRST-SKELETON. Over the PID \(\mathbf Q[t]\), the flatness
hypothesis is equivalent to absence of \(t\)-torsion in these
graded modules.

\emph{Proof and bounds.}
The vertical cohomology of each factor is its coefficient copy
in degree zero plus modules supported in bounded weights; the
coefficient copy splits by a vertex map. Choose bounds
\(a_i,b_i\geq1\) for those defect modules, uniform in the simplex.
The graded K\"unneth exact sequence over the PID
\(\mathbf Q[t]\) contains tensor terms with \(i+j=q\) and
\(\operatorname{Tor}_1\) terms with \(i+j=q+1\).
It applies here to the nonnegative flat complexes; in a fixed
output degree only finitely many terms are involved.

If modules are supported in weights less than \(a,b\), their
tensor is supported in weights less than \(a+b\), and so is
their \(\operatorname{Tor}_1\). For the latter assertion use
the canonical free resolution of a \(\mathbf Q[t]\)-module \(M\):
its two free terms are
\(\mathbf Q[t]\otimes_{\mathbf Q}M\), with one shifted upward
by weight one, and differential \(t\otimes m-1\otimes tm\).
After tensoring with \(N\), the kernel is supported in weights
at most \((a-1)+(b-1)+1\).
Terms involving a coefficient copy of \(\mathbf Q[t]\) produce
no Tor. A safe vertical cutoff in output degree \(q\) is
\[
                  \max_{i+j\leq q+1}(a_i+b_j).
 \tag{136.LT3}
\]

For a form degree \(i\), let \(\rho_i\) bound the weight support
of all horizontal homology of the \(C\)-row, and let \(\sigma_j\)
do the same for \(D\). Normalized modules are direct summands
of the simplex modules and remain flat. The
Eilenberg--Zilber comparison followed by the same K\"unneth
calculation gives the horizontal cutoff
\[
                  \max_{i+j=q}(\rho_i+\sigma_j).
 \tag{136.LT4}
\]
The bounds do not depend on horizontal degree or simplex.
They prove the asserted rational admissibility.
\hfill\(\square\)

\paragraph{ARCH-136-NONFLAT-TENSOR.}
The flatness assumption in the preceding statement cannot be
discarded, even for theories with the same standard first
skeleton and with exact vertical coefficient cohomology in
every weight.

\emph{Construction.}
Let \(B\) be the inherited standard integral PD simplex model.
In each simplex, changing coordinates to
\(t,x_1,\ldots,x_n\) makes \(B_n\) a tensor of free even-to-odd
PD disks over \(S\). In particular its terms are \(S\)-free,
and it contracts vertically to \(S\).
For every \(m\geq3\) let \(K(m)\) be the exact graded
\(S\)-complex
\[
 S e_m\xrightarrow{\ e_m\mapsto t f_m\ }
 S f_m\xrightarrow{\ f_m\mapsto g_m\ }
 (S/(t))g_m,
 \tag{136.LT5}
\]
in form degrees \(1,2,3\), with
\[
 \operatorname{wt}(e_m)=m+1,\qquad
 \operatorname{wt}(f_m)=\operatorname{wt}(g_m)=m.
 \]
Multiplication by \(t\) is injective in \(S\), so exactness is
literal, including degree three.
Let \(K=\bigoplus_{m\geq3}K(m)\), and let
\[
 U=DK(\mathbf Z[3]\xrightarrow{\,1\,}\mathbf Z[2]).
 \]
Then \(U\) is simplicially contractible,
\(U_0=U_1=0\), and \(U_2=\mathbf Z\). Put
\[
 \mathcal I_n=(B_n\otimes_S K)\otimes_{\mathbf Z}U_n,
                 \qquad C=B\oplus\mathcal I,
 \tag{136.LT6}
\]
with square-zero multiplication on \(\mathcal I\).
Use the split PD structure: all divided powers of order at
least two on an even element of \(\mathcal I\) are zero, and
\[
 \gamma_r(b+a)=\gamma_r(b)+\gamma_{r-1}(b)a
 \quad(r\geq2)
 \tag{136.LT7}
\]
for a homogeneous even \(a\in\mathcal I\).
This is the ordinary split square-zero PD structure. The
special differential identity follows by differentiating
(136.LT7); products of two ideal elements vanish.

Each \(\mathcal I_n\) is vertically acyclic: the terms of \(B_n\)
are \(S\)-free, its form degrees are bounded, and \(U_n\) is
free abelian. Each horizontal row of \(\mathcal I\) is
contractible, because the simplicial contraction of \(U\)
tensors with the degeneracies of \(B\). Consequently \(C\)
has exactly the vertical and horizontal cohomology of \(B\),
with its uniform cutoffs. It is admissible, and
\[
                         C_0=S,\qquad C_1=B_1.
 \]

Now rationalize and write \(R'=\mathbf Q[t]\). The ordinary
tensor \(C\otimes_SC\) has the direct summand
\(\mathcal I\otimes_S\mathcal I\) as a complex. In simplex two
its rationalization is
\[
 (B_{2,\mathbf Q}\otimes_{R'}B_{2,\mathbf Q})
       \otimes_{R'}(K_{\mathbf Q}\otimes_{R'}K_{\mathbf Q}).
 \]
The first factor has a weight-preserving \(R'\)-linear
de Rham contraction to \(R'\). Tensoring that explicit
contraction is valid even with nonflat modules. Thus the
vertical cohomology is that of \(K_{\mathbf Q}\otimes_{R'}K_{\mathbf Q}\).

For every pair \(m,n\),
\[
 H^4(K(m)_{\mathbf Q}\otimes_{R'}K(n)_{\mathbf Q})
           \cong R'/(t)
 \quad\hbox{in weight }m+n+1.
 \tag{136.LT8}
\]
Here is the direct differential check, suppressing indices:
\[
 \begin{split}
 d(e\otimes f)&=t(f\otimes f)-e\otimes g,\\
 d(f\otimes e)&=g\otimes e+t(f\otimes f),\\
 d(f\otimes f)&=g\otimes f+f\otimes g,\\
 d(e\otimes g)&=d(g\otimes e)=0.
 \end{split}
 \]
The last two terms are annihilated by \(t\).
A degree-four cycle has its \(f\otimes f\) coefficient in
\(tR'\); modulo the first two boundaries it is represented
by \(e\otimes g=-g\otimes e=t(f\otimes f)\), and the only
remaining relation is its annihilation by \(t\).
This proves (136.LT8). As \(m,n\) grow, it gives nonzero
vertical cohomology in the fixed form degree four in
arbitrarily large weights. The ordinary rational tensor is
therefore not admissible.
\hfill\(\square\)

This example is compatible with the separate positive
standard-factor route: \(B_n\) itself has an \(S\)-linear
vertical contraction, so tensoring an arbitrary \(C\) with
that particular factor preserves its vertical cohomology.
It does not supply the corresponding integral horizontal
acyclicity.

\paragraph{ARCH-136-TORSION-QUOTIENT-PERIOD.}
Let \(A\) be a nonnegative simplicial differential PD
\(S\)-algebra with \(A_0=S\). Suppose \(A[t^{-1}]\) has
vertical coefficient cohomology \(\mathcal R\) in degree
zero and zero in positive degrees. Put
\[
 J=\ker(A\longrightarrow A[t^{-1}]),\qquad A^{\mathrm{tf}}=A/J.
 \tag{136.LT9}
\]
Then \(J\) is a simplicial differential PD ideal,
\(A^{\mathrm{tf}}_0=S\), and every simplex of
\(A^{\mathrm{tf}}\) has \(H^0=S\) exactly.

For a weight \(w\geq0\) and \(a\in A_1^{w,1}\) with \(da\in J\),
choose \(u\in A_1[t^{-1}]^{w,0}\) with \(du=a\).
The endpoint difference
\[
       \operatorname{per}_w(a)=\epsilon_1(u)-\epsilon_0(u)
                      \in\mathcal R^w
 \tag{136.LT10}
\]
is independent of \(u\). It descends to a well-defined map
\[
 H^1((A^{\mathrm{tf}}_1)^{w,*})
                  \longrightarrow \mathcal R^w/S^w.
 \tag{136.LT11}
\]
Its vanishing is exactly the strengthened period condition
that all such \(a\), including those only closed after
localization, have integral periods.

\emph{Proof.}
The kernel \(J\) consists of elements annihilated by some
power of \(t\); it is plainly a differential simplicial
ideal. If \(t^Na=0\), then
\(t^{Nr}\gamma_r(a)=\gamma_r(t^Na)=0\), so it is also a
PD ideal. The quotient embeds into its localization.
If a zero-form in this quotient is closed, its image in
the localization is a coefficient in \(\mathcal R\).
Its value at a vertex belongs to \(S\). Injectivity into
the localization then identifies the original zero-form
with that coefficient of \(S\).

Two choices of \(u\) differ by a degree-zero closed
element of the localization, hence by a constant
coefficient. This proves independence. Changing \(a\)
by \(J\) changes nothing after localization, while a
boundary \(db\) has period
\(\epsilon_1(b)-\epsilon_0(b)\in S^w\). This proves
(136.LT11) and its exact stated criterion.
\hfill\(\square\)

If \(A^{\mathrm{tf}}\) maps to an admissible theory \(E\)
with \(E_0=S\), then (136.LT11) must vanish in every weight
where \(H^1(E_1^{w,*})=0\). Indeed an integral primitive in
\(E_1\) has integral endpoint difference, and after
localization periods are functorial. This is a necessary
condition for that quotient-based completion. It is not
claimed for noncycles in \(A\) before the torsion quotient.

For the common first skeleton \(F_1\), multiplication
\(F_1\otimes_SF_1\to F_1\) proves the strengthened condition
for the initial tensor. The target \(F_1\) is \(S\)-free
and vertically exact in every weight. Thus an element
whose differential is \(t\)-torsion maps to an actual
closed element, with an integral primitive in \(F_1\).
Multiplication agrees with the identity on the vertex
coefficient algebra, so its period equals the original
localized period. For the standard-factor tensor the direction
is instead \(F_1\to B_1\). In the inherited first skeleton
\[
 F_1=S\langle x_m:m\geq N\rangle_{\mathrm{PD}}
                     \otimes\Lambda(e_m:m\geq N),
 \qquad dx_m=e_m,
\]
the cutoff \(N\) need not be one. The explicit
endpoint-preserving strict PD map is
\[
 x_m\longmapsto\gamma_m(x),\qquad
 e_m\longmapsto\gamma_{m-1}(x)e
 \quad\hbox{in }B_1=S\langle x\rangle_{\mathrm{PD}}\otimes\Lambda(e).
\]
Its endpoint values are \(0,T^{[m]}\), as required.
Tensoring this map with \(B_1\) and then multiplying gives
\(B_1\otimes_SF_1\to B_1\). The target \(B_1\) is
\(S\)-free and vertically exact, so the preceding period
argument applies without a weight-one generator in \(F_1\).

\paragraph{ARCH-136-PERIOD-REPAIR-FAILURES.}
Neither a localized retraction nor the separate integrality
of the endpoint values of a newly adjoined zero-form
automatically preserves this period condition.

First choose \(h\in B_1^{w+N,0}\), \(w,N\geq1\), with
endpoint difference \(T^{[w+N]}\). Adjoin a reduced pair
\[
 dc=0,\qquad dt'=t^N c-dh,
 \quad |c|=1,\quad |t'|=0,\quad
 \operatorname{wt}(c)=w,\quad\operatorname{wt}(t')=w+N,
 \tag{136.LT12}
\]
with both endpoint values of \(t'\) equal to zero.
This is the zero-cycle instance of a proposed coefficient-power
Tate repair. It has a localized retraction
\[
                  c\longmapsto t^{-N}dh,\qquad t'\longmapsto0.
 \]
After localization it is a contractible even-to-odd disk.
Nevertheless the localized primitive \((t'+h)/t^N\) of
the integral closed element \(c\) has period
\[
 \frac{T^{[w+N]}}{t^N}
       =\frac{w!}{(w+N)!}\,T^{[w]}\notin S^w.
 \tag{136.LT13}
\]
The pair can be attached simplicially on
\(\mathbf Z[\Delta[1]]/\mathbf Z[\partial\Delta[1]]\);
all differentials respect the specified vertices.
This is a bad, avoidable choice of a repair, not a
nonexistence theorem for coherent choices.

Second, write the standard edge as
\[
 B_1=\Gamma_{\mathbf Z}(T,x)\otimes\Lambda(e),
       \qquad dx=e,\qquad \epsilon_0(x)=0,\quad\epsilon_1(x)=T.
 \]
For any prime \(p\), adjoin inside its rationalization the
zero-form
\[
                         f=\frac{x(T-x)}p.
 \]
It has weight two and both endpoint values zero. But \(fe\)
is an integral element of the enlarged algebra, is closed,
and its rational primitive has period
\[
            \int_0^T\frac{x(T-x)}p\,dx
                         =\frac{T^3}{6p}
                         =\frac1p\,T^{[3]}.
 \tag{136.LT14}
\]
The simplicial PD subalgebra generated by \(B\) and this
element retains the vertex coefficient algebra \(S\),
but fails the period condition. Moreover
\(T^{[w-3]}fe\) has period
\(\binom w3 T^{[w]}/p\). For
\(w=3+p^r,\ p^r>3\), that binomial coefficient is one
modulo \(p\). The defect therefore occurs in arbitrarily
large weights in the fixed form degree one.
This is a concrete instance of the separate period-lattice
rigidity result; it shows why endpoint compatibility of
one primitive alone cannot justify an iterative repair.
\hfill\(\square\)

\paragraph{Remaining gate and verification scope.}
FIRST-SKELETON gives incoming admissible replacements
\(C'\to C\) and \(D'\to D\), with shared free first
skeleton. A common admissible target of the replacements
would indeed give connectedness via this zigzag. The
results above leave the coherent integral period lattice
through the higher repairs unresolved. They do not assert
that the arbitrary choices in (136.LT12) or (136.LT14)
are forced.
The complete saved module was read by the author, root, and
the independent reviewer \texttt{/root/cote1348}, with PASS
in these exact scopes (GPT-6 Astra Pro; effort not exposed;
2026-11-22). A later dependency recheck by that reviewer
corrected the direction of the standard-factor edge map:
it is \(F_1\to B_1\), not a map requiring a weight-one disk
in \(F_1\). The earlier author and peer readings had missed
this point; the corrected argument is explicit above and
was reread. The register distinguishes these saved-file
readings from the proofs communicated during the temporary
execution outage.
